\[ %%%%%%%%%%%%%%%%%%%%%%%%%%%%%%%%%%%%%%%%%%%%%%%%%%%%%%%%%%%%%%%%%%%%%%%%%%%%%%%% \subsection{Latent Physical State Evolution} \label{sec:latent_dynamics} %%%%%%%%%%%%%%%%%%%%%%%%%%%%%%%%%%%%%%%%%%%%%%%%%%%%%%%%%%%%%%%%%%%%%%%%%%%%%%%% The latent physical state evolution module defines the temporal dynamics of PHYSGEN by modeling the evolution of hidden representations that encode the essential physical structure of the system. The fundamental assumption is that complex physical processes may contain high-dimensional observable states while evolving according to lower-dimensional latent dynamics governed by underlying physical principles. Instead of directly learning a mapping between consecutive observations: \begin{equation} X_t \rightarrow X_{t+1}, \end{equation} PHYSGEN introduces an intermediate latent physical state: \begin{equation} H_t=f_{enc}(X_t), \end{equation} where $H_t$ represents a physically meaningful latent representation generated by the encoder. The temporal evolution of the system is then modeled as: \begin{equation} H_{t+\Delta t} = \Psi_{\theta} ( H_t,G_t,P_t ), \end{equation} where: \begin{itemize} \item $H_t$ represents the current latent physical state; \item $G_t$ represents the dynamic interaction graph; \item $P_t$ represents physical parameters; \item $\Psi_{\theta}$ denotes the learnable latent dynamics operator. \end{itemize} %%%%%%%%%%%%%%%%%%%%%%%%%%%%%%%%%%%%%%%%%%%%%%%%%%%%%%%%%%%%%%%%%%%%%%%%%%%%%%%% \subsubsection{Latent Dynamics as a Physical Evolution Operator} %%%%%%%%%%%%%%%%%%%%%%%%%%%%%%%%%%%%%%%%%%%%%%%%%%%%%%%%%%%%%%%%%%%%%%%%%%%%%%%% The latent evolution operator is designed as a neural approximation of the underlying physical transition process: \begin{equation} \Psi_{\theta} \approx \mathcal{T}_{phys}, \end{equation} where $\mathcal{T}_{phys}$ represents the unknown physical evolution operator. The continuous-time formulation can be expressed as: \begin{equation} \frac{dH(t)}{dt} = f_{\theta} ( H(t),G(t),P(t) ), \end{equation} which connects PHYSGEN with neural differential equation approaches. However, unlike conventional Neural ODE formulations, where the dynamics function is typically learned without explicit structural constraints, PHYSGEN incorporates physical guidance into the latent evolution process. %%%%%%%%%%%%%%%%%%%%%%%%%%%%%%%%%%%%%%%%%%%%%%%%%%%%%%%%%%%%%%%%%%%%%%%%%%%%%%%% \subsubsection{Physics-Guided Latent Transition} %%%%%%%%%%%%%%%%%%%%%%%%%%%%%%%%%%%%%%%%%%%%%%%%%%%%%%%%%%%%%%%%%%%%%%%%%%%%%%%% The latent transition is decomposed into a learned component and a physical correction component: \begin{equation} H_{t+\Delta t} = H_t + \Delta t ( f_{\theta}(H_t,G_t) + f_{phys}(H_t,P_t) ), \end{equation} where: \begin{itemize} \item $f_{\theta}$ represents learned nonlinear dynamics; \item $f_{phys}$ represents physics-guided corrections. \end{itemize} This formulation allows the model to combine flexible data-driven approximation with explicit physical reasoning. The physical correction term may encode: \begin{itemize} \item conservation constraints; \item known differential operators; \item constitutive relationships; \item symmetry properties; \item domain-specific physical priors. \end{itemize} %%%%%%%%%%%%%%%%%%%%%%%%%%%%%%%%%%%%%%%%%%%%%%%%%%%%%%%%%%%%%%%%%%%%%%%%%%%%%%%% \subsubsection{Latent Physical Manifold} %%%%%%%%%%%%%%%%%%%%%%%%%%%%%%%%%%%%%%%%%%%%%%%%%%%%%%%%%%%%%%%%%%%%%%%%%%%%%%%% PHYSGEN assumes that physically valid system states occupy a constrained region of the latent space, referred to as the latent physical manifold: \begin{equation} \mathcal{M}_{phys} \subset \mathbb{R}^{d}. \end{equation} The objective of latent dynamics evolution is therefore not only to predict future states but also to maintain trajectories within this physically admissible manifold: \begin{equation} H_t \in \mathcal{M}_{phys} \quad \forall t. \end{equation} To enforce this property, PHYSGEN introduces a manifold consistency functional: \begin{equation} \mathcal{L}_{manifold} = d ( H_t, \mathcal{M}_{phys} ), \end{equation} where $d(\cdot)$ measures deviation from the physically valid latent region. %%%%%%%%%%%%%%%%%%%%%%%%%%%%%%%%%%%%%%%%%%%%%%%%%%%%%%%%%%%%%%%%%%%%%%%%%%%%%%%% \subsubsection{Multi-Step Latent Rollout} %%%%%%%%%%%%%%%%%%%%%%%%%%%%%%%%%%%%%%%%%%%%%%%%%%%%%%%%%%%%%%%%%%%%%%%%%%%%%%%% Long-horizon forecasting requires repeated application of the latent evolution operator: \begin{equation} H_{t+k} = \Psi_{\theta}^{k} ( H_t,G_t ). \end{equation} However, recursive prediction introduces error accumulation: \begin{equation} \epsilon_{t+k} = || H_{t+k} - \hat{H}_{t+k} ||. \end{equation} PHYSGEN reduces this accumulation by combining: \begin{itemize} \item physics-guided latent corrections; \item stability-aware state regulation; \item invariant preservation mechanisms. \end{itemize} %%%%%%%%%%%%%%%%%%%%%%%%%%%%%%%%%%%%%%%%%%%%%%%%%%%%%%%%%%%%%%%%%%%%%%%%%%%%%%%% \subsubsection{Connection with Neural Operators} %%%%%%%%%%%%%%%%%%%%%%%%%%%%%%%%%%%%%%%%%%%%%%%%%%%%%%%%%%%%%%%%%%%%%%%%%%%%%%%% The latent evolution module can also be interpreted as a graph-based neural operator: \begin{equation} \Psi_{\theta}: (H_t,G_t) \rightarrow H_{t+\Delta t}. \end{equation} Unlike classical neural operators that primarily learn mappings between function spaces, PHYSGEN learns a physically guided evolution process over dynamically changing interaction structures. This enables the framework to combine: \begin{itemize} \item operator learning capability; \item graph-based relational reasoning; \item physical consistency constraints. \end{itemize} %%%%%%%%%%%%%%%%%%%%%%%%%%%%%%%%%%%%%%%%%%%%%%%%%%%%%%%%%%%%%%%%%%%%%%%%%%%%%%%% \subsubsection{Summary} %%%%%%%%%%%%%%%%%%%%%%%%%%%%%%%%%%%%%%%%%%%%%%%%%%%%%%%%%%%%%%%%%%%%%%%%%%%%%%%% The latent physical state evolution module provides the temporal foundation of PHYSGEN. By modeling system dynamics in a constrained latent physical space, the framework aims to achieve stable long-horizon forecasting while preserving the structural properties of the underlying physical process. Together with physics-guided message passing, this module transforms PHYSGEN from a predictive graph neural network into a physics-guided dynamical system model capable of continuous physical evolution. \] 